In the \VH{} analysis, at least one primary vertex is required to be reconstructed each with at least two associated tracks with $\pt > 500$ MeV. If multiple vertices are reconstructed, the one with the largest $\sum\limits_{\mathrm{Tracks}} \pt^{2}$ is selected as the primary vertex. 

Leptons in this analysis are required to originate from the primary vertex. To ensure this, the absolute value of the longitudinal impact parameter $z_0$ projected onto the beam line with $\sin{\theta}$ ($|\Delta{z_0}\sin{\theta}|$) must be less than 0.5 mm. Additionally, the significance of the transverse impact parameter calculated with respect to the beam line ($\frac{|d_0|}{\sigma_{d_{0}}}$) must be less than three for muons and five for electrons.

Electrons in the 3$\ell$ channel use the ``TightLH'' identification and ``Tight\_VarRad'' isolation WP's while electrons in the 4$\ell$ use the ``MediumLH'' identification and ``Loose\_VarRad'' WP's\@. Different from previous analyses, electrons that are reconstructed as ambiguous (e.g.\ cannot be distinguished from a photon) are kept in the analysis, as well as electrons in the transition region ($1.37 < |\eta| < 1.52$), extending electron coverage to the full $|\eta| < 2.47$.

Muons in the 3$\ell$ and 4$\ell$ channels are required to pass the `Medium' quality criteria, while the 3$\ell$ requires the `Tight\_VarRad' isolation WP and the 4$\ell$ requires the `Loose\_VarRad' WP\@. Additionally, muon candidates are required to have $\pt > 10$ GeV, and fall in the range $|\eta| < 2.5$. 

Jets in this analysis are required to have $\pt > 30$ GeV and $|\eta| < 4.5$, however, subthreshold jets with $20 < \pt < 30$ are included for the purpose of $b$-tagged jet counting. For jets with \pt{} between 20 and 60 GeV and within $|\eta| < 2.5$, they are required to pass NNJvt with the `FixedEffPt' WP, while jets in the same \pt{} range but within $2.5 < |\eta| < 4.5$ must pass the fJvt tagger with the `Tighter' WP\@. Jets originating from $b$-quarks with $\pt > 20$ GeV and $|\eta| < 2.5$ are identified using the GN2 $b$-tagging algorithm with the `Continuous' WP with a score of at least two, corresponding to an 85\% $b$-tagging efficiency. 

In this analysis, the $W$ bosons are required to decay leptonically, producing a momentum imbalance in the transverse plane. The \met{} is calculated using the `Tight' WP\@. However, a separate set of leptons are used in this calculation than defined above. Specifically, the input muons must pass the `Loose' quality WP with no isolation requirement and $\frac{|d_0|}{\sigma_{d_{0}}} < 15$, while electrons are required to pass the `LooseLH' identification WP and `Loose\_VarRad' isolation WP\@. Additionally, all jets with $\pt > 20$ GeV are used as input for the \met{} calculation.

To ensure no objects overlap with one another, the following overlap removal criteria is applied:
\begin{itemize}
  \item Electron-Electron: If there is an overlapping second layer cluster or a shared track, the lower \pt{} electron is removed.
  \item Electron-Muon: If a combined muon shares an ID track with an electron, the electron is removed. If a calorimeter tagged muon shares an ID track with an electron, the muon is removed.
  \item Electron-Jet: If $\Delta{R}(\mathrm{Jet}, e) < 0.2$ the jet is removed. For all other jets, if $\Delta{R}(\mathrm{Jet}, e) < \min(0.4, 0.04 + 10\mathrm{GeV}/\pt^{e})$ the electron is removed.
  \item Muon-Jet: If $\Delta{R}(\mathrm{Jet}, e) < 0.2$ the jet is removed if it has less than three tracks with $\pt > 500$~MeV. If $\pt^{\mu}/\pt^{\mathrm{Jet}} > 0.5$ and $\pt^{\mu}/\sum\limits_{\mathrm{Tracks}}\pt^{\mathrm{Jet}} > 0.7$ the jet is also removed. For any leftover jets, if $\Delta{R}(\mathrm{Jet}, \mu) < \min(0.4, 0.04 + 10\mathrm{GeV}/\pt^{\mu})$ the muon is removed.
\end{itemize}

The overlap removal process is $b$-tagged jet aware, meaning that if a jet is $b$-tagged, it will never be removed. Object pre-selections used for overlap removal are listed in Table~\ref{tab:vh_objects_overlap_removal}

\begin{table}[ht]
\centering
\begin{tabular}{c|c}
  \hline
  \multicolumn{2}{c}{\textbf{Electrons}} \\ 
  \hline
  \pt{} & $> 10$ GeV \\ 
  $\eta$ & $|\eta| < 2.5$ \\ 
  ID & \texttt{VeryLooseLH} \\
  Isolation & None \\
  Impact Parameters & $|z_{0}\sin{\theta}| < 0.5$ mm, $|d_{0}/\sigma_{d_{0}}| < 5$ \\
  \hline
  \multicolumn{2}{c}{\textbf{Muons}} \\ 
  \hline
  \pt{} & $> 10$ GeV \\
  $\eta$ & $|\eta| < 2.5$ \\
  Quality & \texttt{Loose} \\
  Isolation & None \\
  Impact Parameters & $|z_{0}\sin{\theta}| < 0.5$ mm, $|d_{0}/\sigma_{d_{0}}| < 15$ \\
  \hline
  \multicolumn{2}{c}{\textbf{Jets and $b$-Jets}} \\
  \hline
  \pt{} & $> 20$ GeV \\
  $\eta$ & $|\eta| < 4.5$ \\
  JVT, FJVT & \texttt{FixedEffPt}, \texttt{Tighter} \\
  $b$-jet & \texttt{GN2} with 85\% WP \\
  \hline
\end{tabular}
\caption{Object selection used in overlap removal}\label{tab:vh_objects_overlap_removal}
\end{table}